De oude unix systemen waren voorzien van terminals\index{terminal}. Een terminal is niets anders dan een toetsenbord en een scherm. De toetsaanslagen gingen digitaal over de lijn naar een server en deze verwerkte de dat, een kopie werd naar het scherm gestuurd zodat de gebruiker ook wist wat deze getypt had. Deze oplossing was feitelijk een ouderwetse typemachine over een electronisch lijntje, het kreeg dan ook de naam teletypewriter\index{teletypewriter} (afgekort TTY\index{TTY}). Tele van ver en typewriter van typemachine.

Met de komst van netwerken was het onhandig om en allemaal lijntjes te leggen van de server naar de terminals en netwerkkabels te leggen naar de computers. De terminals kregen ook een netwerk aansluiten en een nieuw protocol werd bedacht om de data van deze netwerkterminals\index{netwerkterminal} naar de server te krijgen. Dat protocol werd \textbf{telnet}\index{telnet} genoemd (port 23). Alle data van deze systemen ging voornamelijk over lokale netwerken en dus ging alle data in leesbare tekst over de lijnen. Telnet werd ook veel toegepast om routers en andere netwerkapparatuur te configureren.

Over het Internet was het gebruik van telnet geen veilig protocol. Bij het verbinden aan routers en servers moesten gebruikers gebruikersnamen en wachtwoorden ingeven en ook deze gingen in leesbare tekst over de netwerkverbindingen. Met de komst van netwerksniffers werd dat een steeds groter veiligheidsrisico. Een nieuwe protocol was nodig om te zorgen dat de data veilig over de lijn zou gaan. Dit protocol is Secure SHell\index{Secure Shell} ofwel SSH\index{SSH}.

